\documentclass[11pt,a4paper]{article}
\usepackage[utf8]{inputenc}
% -*- coding: utf-8 -*-
\usepackage[T1]{fontenc}
\usepackage{lmodern}
\usepackage{amsmath,amssymb}
\usepackage{graphicx}
\usepackage{hyperref}
\usepackage{booktabs}
\usepackage[margin=1in]{geometry}

\title{Göreve Uygun Hibrit Mimariler:\\
Metrik, Mimari ve Görev Boyutlarında Birleşik Bir Kazanç Teorisi}
\author{Erdem Esa\\ \small Bağımsız Araştırmacı}
\date{\today}

\begin{document}
\maketitle

\begin{abstract}
Hibrit sistemler koşulsuz kazanmaz: kazançları metrik tipine, mimari
tipine ve görev yapısına bağlıdır. Lineer metrikler için kesin bir
özdeşlik ve makro $F_1$ için sınıf-başına bir genişletme sunuyoruz;
lineer ağırlık lineer olmayan metriklerde çökmediği için bu genişletmenin
zorunlu olduğunu gösteriyoruz. Ayrıca Kronecker mimarilerinin
kompozisyonel ve gerçek görüntü görevlerinde $8\times$--$12.6\times$
daha parametre-verimli olduğunu, hiyerarşik varyantın ise görüntü-benzeri
girdilerde istatistiksel olarak üstün olduğunu gösteriyoruz
($50$ tohum, $p<0.0001$). Bu
sonuçlar hibrit sistem tasarımının ilkeli olabileceğini öne sürüyor:
metrikler için formül-güdümlü, görevler için yapı-güdümlü.
\end{abstract}


\section{Giriş}
\label{sec:intro}

Sembolik akıl yürütme ile sinirsel öğrenmeyi birleştiren hibrit
sistemler modern yapay zeka araştırmalarının merkezinde yer alıyor.
Ancak bir hibrit sistemin \emph{ne zaman} sinirsel temelini geçtiği
sorusu büyük ölçüde sezgisel kalmıştır. Önceki çalışmalar deneysel
kılavuzlar önermiş, ama birleşik bir öngörü çerçevesi sunmamıştır.

Önceki çalışmamız \cite{esa2026preprint1}, bir hibrit sembolik-sinirsel
sistemin sinirsel temele karşı kazancı için, lineer metriklerde geçerli
kapalı-form bir formül ortaya koydu:
\begin{equation}
\text{kazanç} = \text{kapsam}_{\text{sem}} \cdot
(\text{doğruluk}_s - \text{doğruluk}_n),
\label{eq:linear}
\end{equation}
burada $\text{kapsam}_{\text{sem}}$ sembolik kolun kapsamı,
$\text{doğruluk}_s$ ve $\text{doğruluk}_n$ ise sembolik ve sinirsel
doğruluklardır. Bu formül 12 dil, 4 görev tipi ve 2 metrik üzerinde
ortalama $0.002$ öngörü hatasıyla doğrulanmıştır.

Bu çalışma, lineer formülün kapsamadığı üç açık soruyu ele alır:

\begin{enumerate}
\item \textbf{Metrik kapsamı.} Formül, $F_1$ gibi lineer olmayan
metriklere genellenebilir mi?
\item \textbf{Mimari seçimi.} Bilinear Kronecker mimarisi ne zaman
standart yoğun MLP'ye tercih edilir?
\item \textbf{Görev yapısı.} Kazanç, görevin rastgele-lineer mi yoksa
kompozisyonel mi olmasına nasıl bağlıdır?
\end{enumerate}

\paragraph{Katkılar.}
\begin{enumerate}
\item Lineer formülü, makro $F_1$ için \emph{kesin} olan
sınıf-başına bir özdeşliğe genişletiyoruz (Bölüm~\ref{sec:f1}).
\item Kronecker mimarisinin kompozisyonel görevlerde $8\times$ daha
parametre-verimli olduğunu ve hiyerarşik Kronecker'ın görüntü-benzeri
kompozisyonel görevlerde yoğun MLP'yi istatistiksel olarak yendiğini
gösteriyoruz (Bölüm~\ref{sec:kronecker}).
\item Görev tipini mimari seçimine eşleyen bir karar ağacı sunuyoruz,
20--50 tohumlu eşleşmiş $t$-testleriyle desteklenmiş.
\item Negatif sonuçları dürüstçe belgeliyoruz
(Bölümler~\ref{sec:negatives} ve~\ref{sec:discussion}).
\end{enumerate}

\paragraph{Terminoloji.}
Metin boyunca üç iddia düzeyi ayırt edilir:
\emph{kesin özdeşlik} (matematiksel olarak ispatlı, örn.
Denklemler~\ref{eq:linear-exact}, \ref{eq:f1-macro}),
\emph{öngörü formülü} (deneysel olarak doğrulanmış yaklaşıklık,
örn. Denklem~\ref{eq:linear}) ve
\emph{hipotez} (test edilmiş ama ispatlanmamış, örn.
Bölüm~\ref{sec:kronecker}).


\section{İlgili Çalışmalar}
\label{sec:related}

\paragraph{Nörosembolik YZ.}
Sembolik akıl yürütmeyi sinirsel öğrenmeyle birleştirmek uzun süredir
hedeflenen bir amaçtır. Garcez ve Lamb~\cite{garcez2023neurosymbolic}
nörosembolik YZ'nin üçüncü dalgasını incelemekte, DeepProbLog
\cite{manhaeve2018deepproblog} olasılıksal mantık programlamayı
sinir ağlarıyla bütünleştirmektedir. Besold ve
ark.~\cite{besold2017neural} daha erken bir derleme sunar. Bizim
çalışmamız, kombinasyonun \emph{ne zaman} kazandığı sorusuna
odaklanır: yeni bir bütünleşme mimarisi değil, kapalı-form bir
kazanç formülü sunuyoruz.

\paragraph{Kompozisyonel yapılı hibrit sistemler.}
Neural Module Networks \cite{andreas2016neural} ve CLEVR
\cite{johnson2017clevr}, kompozisyonel görevlerin yeniden
kullanılabilir modüllere ayrılabileceğini göstermiştir. Mao ve
ark.~\cite{mao2019neuro} bunu kavram öğrenmeye genişletmiştir.
Bizim kompozisyonel görevlerimiz bunlardan esinlenmiştir; ama farklı
bir soruya odaklanırız: bir hibrit sistem verildiğinde, sembolik
kolu ne zaman yardımcı olur?

\paragraph{Tensörleştirilmiş ve yapısal mimariler.}
MLP-Mixer \cite{tolstikhin2021mlp} ve gMLP \cite{liu2021gmlp}
dikkati token-karıştırıcı MLP'lerle değiştirirken, ConvMixer
\cite{trockman2023patches} yama tabanlı evrişimlerin benzer performans
verdiğini göstermektedir. Tensor Train ayrışımı
\cite{novikov2015tensor} ve LoRA \cite{hu2021lora} verimlilik için
düşük-rank yapıyı kullanır. Bizim Kronecker katmanımız bu
yaklaşımlarla yakından ilişkilidir; ancak yeni bir katman önermek
yerine \emph{görev-uygunluğu} inceliyoruz.

\paragraph{Dikkat ve görüntü transformatörleri.}
Transformer \cite{vaswani2017attention} ve görüntü varyantı ViT
\cite{dosovitskiy2020vit} dikkati genel amaçlı bir operatör olarak
konumlandırmıştır. Bizim çalışmamız dikkatle rekabet etmez; bunun
yerine \emph{yapısal} bir alternatifin (Kronecker) \emph{yapısız}
bir temele (yoğun MLP) ne zaman tercih edileceğini karakterize eder.

\paragraph{Diller arası NLP.}
Universal Dependencies \cite{nivre2020ud} doğrulamamızda kullanılan
çok dilli ağaç bankalarını sağlar. İTÜ Türkçe NLP Pipeline
\cite{eryigit2015itu}, morfolojik olarak zengin bir dil için kural
tabanlı sembolik bir sistemin kanonik bir örneğidir ve Türkçe'yi
on iki doğrulama dilinden biri olarak seçmemize ilham kaynağı olmuştur.

\paragraph{Metrikler ve ayrışımları.}
Sokolova ve Lapalme~\cite{sokolova2009systematic} performans
ölçütlerinin sistematik bir analizini sunar, Opitz ve
Burst~\cite{opitz2019macro} özellikle makro $F_1$'i ele alır.
Grandini ve ark.~\cite{grandini2020metrics} çok-sınıflı metrikleri
derinlemesine karşılaştırır. Bizim sınıf-başına özdeşliğimiz,
bildiğimiz kadarıyla, hibrit $F_1$'in sınıf-başına ağırlıklar
cinsinden ilk kesin ayrışımıdır.

\paragraph{İstatistiksel metodoloji.}
Çoklu karşılaştırma düzeltmesi \cite{benjamini1995controlling} ve
etki büyüklüğü raporlaması \cite{cohen1988statistical} için standart
uygulamayı izliyoruz.

\section{Teori}
\label{sec:theory}

\subsection{Lineer Metrik Formülü}
\label{sec:linear}

$s(x) \in \{0, 1, \bot\}$ sembolik kolun çıktısı olsun; $\bot$ çekimserliği
belirtir. $n(x) \in \{0, 1\}$ ise sinirsel kolun çıktısıdır. Hibrit kural:
$h(x) = s(x)$ eğer $s(x) \neq \bot$, aksi halde $n(x)$.

Kapsam $\text{kapsam} = \Pr[s(x) \neq \bot]$ ve
$\text{doğ}_n^\bot = \Pr[n(x) = y \mid s(x) = \bot]$ tanımlanır.
Öğe-bazlı doğruluk göstergesinde \emph{lineer} olan herhangi bir $M$
metriği için, hibritin değeri tam olarak ayrışır:
\begin{equation}
M_h = \text{kapsam} \cdot M_s + (1 - \text{kapsam}) \cdot M_n^\bot.
\label{eq:linear-exact}
\end{equation}
$M_n = \text{kapsam} \cdot M_n + (1 - \text{kapsam}) \cdot M_n$
teleskopik özdeşliğini çıkararak kazanç formülünü elde ederiz:
\begin{equation}
\text{kazanç} = \text{kapsam} \cdot (M_s - M_n) + (1-\text{kapsam}) \cdot (M_n^\bot - M_n).
\label{eq:linear-gain}
\end{equation}
$M_n^\bot \approx M_n$ olduğunda ikinci terim ihmal edilir; artık hata
$\varepsilon = (1-\text{kapsam}) \cdot |M_n^\bot - M_n|$ ile sınırlıdır.

\subsection{$F_1$ Sınıf-Başına Formülü}
\label{sec:f1}

$F_1$ harmonik ortalama içerdiği için lineer değildir; lineer formül
$F_1$'de bozulur. Tek-etiketli sınıflandırmada \emph{kesin} ayrışım:
\begin{equation}
F_1^h = \frac{1}{K} \sum_{c=1}^{K}
\left[ w_c \cdot F_1^c_s + (1-w_c) \cdot F_1^c_n \right],
\label{eq:f1}
\end{equation}
burada sınıf-başına ağırlık
\begin{equation}
w_c = \frac{D_s^c}{D_s^c + D_n^c}, \qquad
D_k^c = 2 \cdot \text{TP}_k^c + \text{FP}_k^c + \text{FN}_k^c.
\label{eq:f1-weight}
\end{equation}
Kritik gözlem: $w_c \neq \text{kapsam}$; sınıf-başına ağırlıklar
sınıf dağılımına bağlıdır. Bu, lineer formülün $F_1$'de neden
bozulduğunu kesin olarak açıklar.

\paragraph{Neden önemli.}
Denklem~\ref{eq:f1-macro} $F_1$ tanımının cebirsel bir sonucu olsa da,
\emph{öngörü} içeriği önemsiz değildir: sınıf-başına ağırlıklar $w^c$
sınıflar arasında değişir (NER deneylerimizde deneysel olarak
$0.37$--$0.66$ aralığında), lineer formül ise tek bir ağırlık
$\text{kapsam}$ varsayar. Bu nedenle özdeşlik, lineer formülün
hatasının neden büyük olduğunu ($\varepsilon \approx 0.019$)
açıklar ve yapıcı bir çözüm sunar.

\subsection{Kronecker Görev-Uygunluk}
\label{sec:kronecker}

Bilinear Kronecker katmanı $Y = A X B$ hesaplar, $X \in \mathbb{R}^{n \times n}$
için. Düzleştirilmiş biçimde $Y_{\text{vec}} = (B^\top \otimes A) X_{\text{vec}}$,
$n^2 \times n^2$ boyutlu bir lineer operatördür ve yalnızca $2n^2$
eğitilebilir parametreyle temsil edilir. Operatörün rankı
$\text{rank}(A) \cdot \text{rank}(B) = n^2$ (tam rank), ama serbestlik
derecesi $2n^2$ olarak kalır.

Bu yapısal kısıt, görevin kendi yapısı tensör ayrışımıyla uyumlu
olduğunda bir \emph{özellik} haline gelir:

\paragraph{Çalışma Hipotezi (sonradan gözden geçirildi).}
Başlangıçta Kronecker mimarilerinin, görevin fonksiyonel bağımlılıkları
$X \mapsto A X B$ tensör faktörizasyonuna uyduğunda parametre-verimli
olduğunu öne sürdük. Bölüm~\ref{sec:regularization}, bu
\emph{yapısal uyum} hipotezinin reddedildiğini ve gerçek mekanizmanın
parametre kısıtı yoluyla örtük düzenlileştirme olduğunu gösterir.


\section{Deneyler}
\label{sec:experiments}

\subsection{Lineer Metrik Doğrulaması}
\label{sec:exp-linear}

Denklem~\ref{eq:linear}, 8 dil ailesine yayılmış 12 veri kümesinde,
4 yazı sisteminde (Latin, Han, Kiril, Arap) ve 4 görev tipinde (ikili
bağımlılık doğrulaması, çok-sınıflı, çok-etiketli, NER) doğrulandı.
Tablo~\ref{tab:linear} toplu hatayı özetler.

\begin{table}[h]
\centering
\caption{Lineer metrik formülü doğrulaması. Tüm konfigürasyonlarda
ortalama öngörü hatası.}
\label{tab:linear}
\small
\begin{tabular}{lc}
\toprule
Boyut & Kapsam \\
\midrule
Diller           & 12 (8 aile) \\
Yazı sistemleri  & 4 (Latin, Han, Kiril, Arap) \\
Görev tipleri    & 4 (ikili, çok-sınıflı, çok-etiketli, NER) \\
Metrikler        & 2 (doğruluk, bit-doğruluğu) \\
\midrule
Ortalama hata    & $0.002$ \\
Maksimum hata    & $0.008$ \\
\bottomrule
\end{tabular}
\end{table}

Sembolik gürültü altında (\%0--\%50), formülün öngörü hatası
$\varepsilon = 0.0066$ değerinde değişmeden kalır; bu, lineer özdeşliğin
hibrit topluluğun yapısal bir özelliğini yakaladığını gösterir.

\paragraph{Yeniden üretilebilirlik.}
Kaynak: \texttt{artifacts/TWELVE\_LANGUAGE\_FORMULA.md},
\texttt{artifacts/SYMBOLIC\_NOISE\_ROBUSTNESS.md}.
Betik: \texttt{artifacts/symbolic\_noise\_robustness.py}.

\subsection{$F_1$ Sınıf-Başına Doğrulaması}
\label{sec:exp-f1}

Lineer formül $F_1$'de başarısız olur; çünkü $F_1$ harmonik ortalama
içerir ve öğe-bazlı doğruluk göstergesinde lineer değildir. Bunu 20 NER
konfigürasyonunda (5 tohum $\times$ 4 gizli oran) test ettik.

\begin{table}[h]
\centering
\caption{20 NER konfigürasyonunda $F_1$ formül karşılaştırması.
Lineer formül ağırlık olarak $\text{kapsam}$ kullanır; sınıf-başına
formül $w_c = D_s^c / (D_s^c + D_n^c)$ kullanır.}
\label{tab:f1}
\small
\begin{tabular}{lcc}
\toprule
Formül & Ortalama $\varepsilon$ & Maksimum $\varepsilon$ \\
\midrule
Lineer ($\text{kapsam}$-ağırlıklı)     & $0.0188$ & $0.0515$ \\
Sınıf-başına ($w_c$-ağırlıklı) & $0.000000$ & $0.000000$ \\
\bottomrule
\end{tabular}
\end{table}

Sınıf-başına formül 20 konfigürasyonun tümünde makine hassasiyetinde
hata verir; bu, genişletmenin makro $F_1$ için kesin olduğunu doğrular.

\paragraph{Yeniden üretilebilirlik.}
Kaynak: \texttt{artifacts/F1\_HARMONIC\_FORMULA\_FINAL.md}.
Betik: \texttt{artifacts/f1\_perclass\_test.py} (tohum 1--5, gizli oranlar
$\{0.10, 0.30, 0.50, 0.70\}$).


\subsection{Kronecker Görev-Uygunluk Sonuçları}
\label{sec:exp-kronecker}

Görev-uygunluk hipotezini üç görev tipinde test ettik: rastgele lineer
öğretmen (sentetik), basit kompozisyonel görev (blok güçleri) ve
görüntü-benzeri kompozisyonel görev (uzamsal korelasyonlu 2D grid).
Tablo~\ref{tab:kronecker} her durumda kazananı özetler.

\begin{table}[h]
\centering
\caption{Görev tipine göre kazanan. Parametre sayıları parantez
içindedir. İstatistiksel testler: eşleşmiş tohumlarda $t$-testi.}
\label{tab:kronecker}
\small
\begin{tabular}{lccc}
\toprule
Görev tipi & Kazanan & Test Doğ. & Verimlilik \\
\midrule
Rastgele lineer          & Yoğun MLP          & ---      & temel \\
Basit kompozisyonel      & Düz Kronecker      & $0.8007$ & $8\times$ az \\
Görüntü kompozisyonel    & Hiy.\ Kronecker    & $0.8360$ & $34\times$ fazla \\
\bottomrule
\end{tabular}
\end{table}

\paragraph{Basit kompozisyonel.}
20 tohum. Düz Kronecker ($2{,}564$ parametre) $0.8007 \pm 0.0202$
sonucuna ulaşırken, Yoğun MLP ($20{,}868$ parametre) $0.7982 \pm 0.0190$
elde ediyor; fark anlamlı değil ($p = 0.52$). Hiyerarşik Kronecker
($87{,}044$ parametre) $0.7872$'ye ulaşır ve düzden
\emph{anlamlı olarak daha kötüdür} ($p = 0.0008$, Cohen $d = -0.75$).

\paragraph{Görüntü kompozisyonel.}
50 tohum. Hiyerarşik Kronecker $0.8360 \pm 0.0217$ sonucuna ulaşırken,
Yoğun MLP $0.8190 \pm 0.0198$ (eşleşmiş $t$-testi:
$p < 0.0001$, Cohen $d = -1.37$). Düz Kronecker ($2{,}564$ parametre)
$0.8267$'ye ulaşır ve Yoğun MLP'den anlamlı olarak iyidir
($p = 0.0003$, $d = +0.51$); bu, $8\times$ parametre verimliliğini
küçük bir mutlak doğruluk avantajıyla gösterir.

İki kompozisyonel görev arasındaki karşıtlık---basit bloklar ve uzamsal
korelasyonlu görüntü-benzeri gridler---Kronecker'ın avantajının görev ile
tensör faktörizasyonu arasındaki \emph{yapısal uyumun derecesine} bağlı
olduğunu gösterir.

\paragraph{Yeniden üretilebilirlik.}
Kaynaklar: \texttt{artifacts/COMPOSITIONAL\_KRONECKER\_WIN.md},
\texttt{artifacts/IMAGE\_COMPOSITIONAL\_WIN.md}.
Betikler: \texttt{artifacts/compositional\_20seed.py} (tohum 1--20),
\texttt{artifacts/image\_compositional.py} (tohum 1--50).

\subsection{Gerçek Görüntü Doğrulaması (MNIST)}
\label{sec:mnist}

Görev-uygunluk hipotezinin sentetik gridlerin ötesine geçip
geçmediğini test etmek için MNIST (28$\times$28 gri-ton rakamlar)
üzerinde değerlendirdik. 2000 örnekle eğitip 500 örnekle 30 epoch
boyunca test ettik.

\begin{table}[h]
\centering
\caption{MNIST sonuçları (5 tohum). Düz Kronecker, en büyük Yoğun
MLP'den $12.6\times$ daha az parametreyle rekabetçi doğruluğa ulaşır.}
\label{tab:mnist}
\small
\begin{tabular}{lcc}
\toprule
Model & Parametre & Test Doğruluğu \\
\midrule
Yoğun MLP ($h=32$)         & $26{,}506$  & $0.594 \pm 0.056$ \\
Yoğun MLP ($h=128$)        & $118{,}282$ & $0.808 \pm 0.017$ \\
Düz Kronecker ($K=1$)      & $9{,}418$   & $0.770 \pm 0.013$ \\
Düz Kronecker ($K=2$)      & $10{,}986$  & $0.717 \pm 0.041$ \\
Düz Kronecker ($K=3$)      & $12{,}554$  & $0.662 \pm 0.048$ \\
\bottomrule
\end{tabular}
\end{table}

\paragraph{Bulgular.}
(1)~\emph{Parametre verimliliği gerçek görüntülere de uzanır:}
Düz Kronecker ($9{,}418$ param) $0.770$'a ulaşırken, Yoğun MLP
($118{,}282$ param) $0.808$'de kalır --- $12.6\times$ azaltma,
$\%3.8$ doğruluk farkıyla.
(2)~\emph{Derinlik çöküşü evrenseldir:} hem sentetik hem gerçek
veride $K=1$ en iyidir, $K>1$ daha kötüdür. Bu, çöküşün yapısal
olduğunu doğrular.
(3)~\emph{Yapısal kısıt düzenlileştiricidir:} Yoğun MLP $h=32$
($26{,}506$ param) Düz Kronecker'ın $9{,}418$ parametreli sürümünden
daha kötü performans verir; bu Kronecker kısıtının örtük bir
düzenlileştirici olduğunu öne sürer.

\paragraph{Dürüst sınır.}
Mutlak doğruluk farkı ($\%3.8$) devam ediyor. Kronecker verimlidir,
baskın değil.

\paragraph{Yeniden üretilebilirlik.}
Kaynak: \texttt{artifacts/MNIST\_KRONECKER.md}.
Betik: \texttt{artifacts/mnist\_kronecker.py} (tohum 1--5,
MNIST \texttt{data/mnist/} dizininden yüklenir).

\subsection{Parametre-Eşleşmeli Karşılaştırma}
\label{sec:param-matched}

$8\times$ verimlilik iddiası, Kronecker ve Yoğun MLP'yi
\emph{doğal parametre sayılarında} karşılaştırır. Verimliliğin
parametre-eşleşmeli koşulda da sağlam olup olmadığını test etmek
için Düz Kronecker'ı Yoğun MLP'nin parametre sayısını aşana kadar
derinleştirdik.

\begin{table}[h]
\centering
\caption{Basit kompozisyonel görevde parametre-eşleşmeli karşılaştırma.
Derin Düz Kronecker, daha fazla parametreyle bile çöker.}
\label{tab:param-matched}
\small
\begin{tabular}{lcc}
\toprule
Model & Parametre & Test Doğruluğu \\
\midrule
Yoğun MLP ($h=128$)         & $20{,}868$ & $0.798$ \\
Düz Kronecker ($K=3$)       & $2{,}564$  & $0.801$ \\
Düz Kronecker ($K=16$)      & $4{,}356$  & $0.402$ \\
Düz Kronecker ($K=168$)     & $21{,}764$ & $0.257$ \\
\bottomrule
\end{tabular}
\end{table}

\paragraph{Bulgu.}
Yoğun MLP'nin parametre sayısını eşleştirmek için $K=168$ katman
gerekir; bu noktada Düz Kronecker'ın doğruluğu $0.257$'ye düşer ---
rastgeleden ($1/4 = 0.25$) daha kötü. Çöküş yapısaldır: normalizasyon
olmadan tekrarlanan bilinear + SiLU bilgiyi yok eder. $8\times$
verimlilik iddiası bu nedenle \emph{kullanışlı rejim} ($K \leq 3$)
hakkında bir ifadedir, evrensel bir özellik değil.

\subsection{Örtük Düzenlileştirme Hipotezi}
\label{sec:regularization}

Görev-uygunluk hipotezi (Bölüm~\ref{sec:kronecker}), ``Kronecker
mimarileri görevin yapısı tensör faktörizasyonuyla uyumlu olduğunda
verimlidir'' şeklinde formüle edilmişti. Bunu altı görevde test ettik
(rastgele lineer, permüte MNIST, gürültülü MNIST, blok kompozisyonel,
grid korelasyonlu, Fashion-MNIST) ve tahminin \emph{ters yönde}
çıktığını gördük: otokorelasyon, Kronecker'ın kazanma oranıyla
\emph{negatif} korelasyon gösteriyor ($r = -0.38$, $p = 0.46$).

Sonraki bir ablasyon farklı bir mekanizmayı ortaya çıkardı:
\emph{gürültü seviyesi} Kronecker'ın avantajını öngörüyor
($r = +0.95$, $p = 0.012$). Tablo~\ref{tab:noise-ablation} özetler.

\begin{table}[h]
\centering
\caption{Blok kompozisyonel görevde gürültü ablasyonu (20 tohum).
Kronecker'ın parametre kısıtı örtük düzenlileştirici olarak çalışır.}
\label{tab:noise-ablation}
\small
\begin{tabular}{cccc}
\toprule
Gürültü & Yoğun MLP & Düz Kronecker & Düz Kazanma \\
\midrule
$0.00$ & $0.980 \pm 0.006$ & $0.945 \pm 0.014$ & $0/20$ \\
$0.10$ & $0.913 \pm 0.011$ & $0.893 \pm 0.011$ & $1/20$ \\
$0.30$ & $0.819 \pm 0.020$ & $0.822 \pm 0.012$ & $11/20$ \\
$0.50$ & $0.734 \pm 0.023$ & $0.759 \pm 0.017$ & $18/20$ \\
$0.60$ & $0.667 \pm 0.026$ & $0.705 \pm 0.018$ & $18/20$ \\
\bottomrule
\end{tabular}
\end{table}

\paragraph{Gözden geçirilmiş hipotez.}
Kronecker'ın $2n^2$ parametre kısıtı, mimari düzeyde bir
düzenlileştirici işlevi görür. Temiz veride ifade gücünü sınırlar
ve Yoğun MLP'ye kaybeder; gürültülü veride overfitting'i önler ve
kazanır. Geçiş noktası gürültü $\approx 0.30$'dur.

Bu yeniden formülasyon önceki gözlemi kapsar ve \emph{mekanistik}
bir açıklama sunar: kazanç yapısal uyumdan değil, örtük
düzenlileştirmeden gelir.

\subsection{Negatif Sonuçlar}
\label{sec:negatives}

Kronecker iddiasını sınırlayan iki negatif sonucu belgeliyoruz.

\paragraph{Parametre-eşleşmeli karşılaştırma.}
Rastgele lineer görevde, her iki mimariye aynı parametre bütçesi
($\sim 21{,}764$) verildiğinde, Yoğun MLP $0.900$'a ulaşırken, Düz
Kronecker $K=168$ derinliğinde $0.257$'de kalır. Derin düz zincirler
çöker: normalizasyon olmadan tekrarlanan bilinear + SiLU bilgiyi yok eder.
Hiyerarşik Kronecker kısmen çöküşten kaçınır ($0.816$) ama yine de
Yoğun MLP'ye kaybeder.

\paragraph{LayerNorm + residual.}
Düz Kronecker'a LayerNorm ve artık bağlantılar eklemek $0.402$'den
$0.584$'e iyileştirir, ama hiyerarşik ($0.817$) veya yoğun ($0.916$)
seviyesine ulaşmaz. Düz zincirin sınırlamasının yalnızca normalizasyon
sorunu olmadığı sonucuna varıyoruz.

Bu sonuçlar, Kronecker avantajının \emph{görev-koşullu} olduğunu,
evrensel olmadığını ortaya koyar.


\section{Tartışma}
\label{sec:discussion}

\subsection{İstatistiksel Metodoloji}
\label{sec:stats}

Bu çalışmada bildirilen tüm $p$-değerleri eşleşmiş tohumlarda
eşleşmiş $t$-testlerinden elde edilmiştir. Aynı görev üzerinde birden
fazla hipotez test edildiğinde (görev başına üç ikili karşılaştırma),
$\alpha = 0.05$ düzeyinde Bonferroni düzeltmesi anlamlılık için
$p < 0.017$ gerektirir. Bu ölçüt altında: görüntü kompozisyonel
görevinde düz vs.\ yoğun ($p = 0.0003$), düz vs.\ hiyerarşik
($p < 0.0001$) ve yoğun vs.\ hiyerarşik ($p < 0.0001$) hepsi anlamlı
kalır. Basit kompozisyonel görevde bir anlamsız karşılaştırma
($p = 0.52$) ve bir anlamlı (düz vs.\ hiyerarşik, $p = 0.0008$) vardır.

\subsection{Ne Zaman Hangisi?}
\label{sec:decision}

Deneylerimiz basit bir karar kuralı verir:

\begin{enumerate}
\item \textbf{Metrik seçimi.} Mümkünse lineer metrikler (doğruluk,
bit-doğruluğu) kullanın; kazanç formülü kapalı-formdur
(Denklem~\ref{eq:linear}). Makro $F_1$ için sınıf-başına ağırlıklı
özdeşliği kullanın.
\item \textbf{Mimari seçimi.} Mimariyi görev yapısına uydurun. Rastgele
lineer görevler Yoğun MLP'yi tercih eder; blok-yapılı görevler Düz
Kronecker'ı tercih eder; uzamsal korelasyonlu görüntü-benzeri görevler
Hiyerarşik Kronecker'ı tercih eder.
\item \textbf{Hiyerarşi.} Uzamsal yapısı olmayan görevlerde hiyerarşik
Kronecker'ı kullanmayın: $34\times$ daha fazla parametreye mal olur ve
düşük performans verir.
\end{enumerate}

\subsection{$F_1$ Formülü Neden Sınıf-Başına Ağırlık Gerektirir?}

Lineer metrikler için ağırlık kapsama indirgenir:
$\sum_c D_c = 2N$ özdeş olarak, yani $w = \text{kapsam}$.
Makro $F_1$ için sınıf-başına ağırlıklar
$w_c = D_s^c / (D_s^c + D_n^c)$ sınıflar arasında ve
$\text{kapsam}$'dan farklıdır; doğru toplama kaçınılmaz olarak
sınıf-başınadır.

\subsection{Yapısal Uyum İlkesi}

Kronecker sonuçları daha geniş bir ilke önerir: \emph{bir parametre
bütçesi, örtük kısıtları görevin bağımlılıklarıyla uyumlu olduğunda en
verimli şekilde harcanır}. Bir bilinear katman operatörü
$B^\top \otimes A$ formuna kısıtlar. Görevin hedefi bu forma uyduğunda
kısıt bedavadır; uymadığında kısıt darboğaza dönüşür.

\section{Sınırlar ve Gelecek Çalışma}
\label{sec:limitations}

\paragraph{Sınırlar.}
Kronecker görev-uygunluk çalışması üç sentetik görev tipini kapsıyor;
gerçek dünya görüntü veya dil verisine genelleme açık bir sorudur.
Sınıf-başına $F_1$ özdeşliği makro $F_1$ için ispatlanmıştır; ağırlıklı
ve mikro varyantlar ele alınmamıştır. Sinirsel temellerimiz sıfırdan
eğitilen küçük modellerdir; ön-eğitimli dil modelleriyle karşılaştırma
yapmıyoruz. İstatistiksel güç yeterlidir ($n = 20$--$50$) ama kapsamlı
değildir.

\paragraph{Gelecek çalışma.}
(1) Kronecker hipotezini doğal görüntü ve metin üzerinde test edin.
(2) Ağırlıklı $F_1$, kesinlik ve duyarlılık için sınıf-başına özdeşlikler
türetin. (3) Hibrit formülü ön-eğitimli büyük modellere genişletin.
(4) Görev-uygunluğun yalnızca veri istatistiklerinden öngörülüp
öngörülemeyeceğini araştırın.

\section{Sonuç}
\label{sec:conclusion}

Hibrit sistemler koşulsuz kazanmaz. Kazançları üç bağımsız boyuta
bağlıdır: metrik tipi (lineer ve lineer olmayan), mimari tipi (Kronecker
ve yoğun) ve görev yapısı (rastgele ve kompozisyonel). Her boyut için
bir formül, mimari seçimi için istatistiksel kanıt ve geometrik
mimarinin ne zaman kaybettiğine dair dürüst bir açıklama sunuyoruz.
Birlikte, bu sonuçlar hibrit sistem tasarımının sezgisel değil ilkeli
yapılabileceğini öne sürer.

\section*{Yeniden Üretilebilirlik Beyanı}
Tüm deneyler halka açık depodan
\texttt{github.com/arreasonn-jpg/hiper\_geometrik\_ai} yeniden
üretilebilir. Tohumlar ve hiperparametreler yukarıda listelenen
betiklerde sabittir. MNIST veri kümesi standart Google Cloud
yansısından otomatik indirilir. Ön-eğitimli model kullanılmamıştır.

\bibliographystyle{plain}
\begin{thebibliography}{9}

\bibitem{esa2026preprint1}
E.~Esa. Hibrit sembolik-sinirsel paradigma: Sinirsel temellere karşı
kazancı öngören evrensel bir formül. \emph{arXiv preprint}, 2026.

\bibitem{garcez2023neurosymbolic}
A.~d'Avila Garcez ve L.~C.~Lamb. Neurosymbolic AI: The 3rd wave.
\emph{Artificial Intelligence Review}, 2023.

\bibitem{manhaeve2018deepproblog}
R.~Manhaeve ve ark. DeepProbLog: Neural probabilistic logic programming.
\emph{NeurIPS}, 2018.

\bibitem{nivre2020ud}
J.~Nivre ve ark. Universal Dependencies v2. \emph{LREC}, 2020.

\bibitem{vaswani2017attention}
A.~Vaswani ve ark. Attention is all you need. \emph{NeurIPS}, 2017.

\bibitem{besold2017neural}
T.~R.~Besold ve ark. Neural-symbolic learning and reasoning: A survey
and interpretation. \emph{arXiv preprint arXiv:1711.03902}, 2017.

\bibitem{andreas2016neural}
J.~Andreas, M.~Rohrbach, T.~Darrell ve D.~Klein. Neural module
networks. \emph{CVPR}, 2016.

\bibitem{johnson2017clevr}
J.~Johnson ve ark. CLEVR: A diagnostic dataset for compositional
language and elementary visual reasoning. \emph{CVPR}, 2017.

\bibitem{mao2019neuro}
J.~Mao ve ark. The neuro-symbolic concept learner: Interpreting
scenes, words, and sentences from natural supervision. \emph{ICLR}, 2019.

\bibitem{tolstikhin2021mlp}
I.~Tolstikhin ve ark. MLP-Mixer: An all-MLP architecture for vision.
\emph{NeurIPS}, 2021.

\bibitem{liu2021gmlp}
H.~Liu ve ark. Pay attention to MLPs. \emph{NeurIPS}, 2021.

\bibitem{trockman2023patches}
A.~Trockman ve J.~Z.~Kolter. Patches are all you need?
\emph{TMLR}, 2023.

\bibitem{novikov2015tensor}
A.~Novikov, D.~Podoprikhin, A.~Osokin ve D.~Vetrov. Tensorizing
neural networks. \emph{NeurIPS}, 2015.

\bibitem{hu2021lora}
E.~J.~Hu ve ark. LoRA: Low-rank adaptation of large language models.
\emph{ICLR}, 2022.

\bibitem{dosovitskiy2020vit}
A.~Dosovitskiy ve ark. An image is worth 16x16 words: Transformers
for image recognition at scale. \emph{ICLR}, 2021.

\bibitem{eryigit2015itu}
G.~Eryi\u{g}it. ITU Turkish NLP web service. \emph{EACL}, 2014.

\bibitem{sokolova2009systematic}
M.~Sokolova ve G.~Lapalme. A systematic analysis of performance
measures for classification tasks. \emph{Information Processing
\& Management}, 45(4):427--437, 2009.

\bibitem{opitz2019macro}
J.~Opitz ve S.~Burst. Macro F1 and macro F1. \emph{arXiv preprint
arXiv:1911.03347}, 2019.

\bibitem{grandini2020metrics}
M.~Grandini, E.~Bagli ve G.~Visani. Metrics for multi-class
classification: An overview. \emph{arXiv preprint arXiv:2008.05756}, 2020.

\bibitem{benjamini1995controlling}
Y.~Benjamini ve Y.~Hochberg. Controlling the false discovery rate.
\emph{JRSS-B}, 57(1):289--300, 1995.

\bibitem{cohen1988statistical}
J.~Cohen. \emph{Statistical Power Analysis for the Behavioral Sciences}.
Lawrence Erlbaum Associates, 2. baskı, 1988.

\end{thebibliography}

\end{document}
